% Options for packages loaded elsewhere
\PassOptionsToPackage{unicode}{hyperref}
\PassOptionsToPackage{hyphens}{url}
\PassOptionsToPackage{dvipsnames,svgnames,x11names}{xcolor}
%
\documentclass[
  11pt,
]{article}
\usepackage{amsmath,amssymb}
\usepackage{iftex}
\ifPDFTeX
  \usepackage[T1]{fontenc}
  \usepackage[utf8]{inputenc}
  \usepackage{textcomp} % provide euro and other symbols
\else % if luatex or xetex
  \usepackage{unicode-math} % this also loads fontspec
  \defaultfontfeatures{Scale=MatchLowercase}
  \defaultfontfeatures[\rmfamily]{Ligatures=TeX,Scale=1}
\fi
\usepackage[]{charter}
\ifPDFTeX\else
  % xetex/luatex font selection
\fi
% Use upquote if available, for straight quotes in verbatim environments
\IfFileExists{upquote.sty}{\usepackage{upquote}}{}
\IfFileExists{microtype.sty}{% use microtype if available
  \usepackage[]{microtype}
  \UseMicrotypeSet[protrusion]{basicmath} % disable protrusion for tt fonts
}{}
\makeatletter
\@ifundefined{KOMAClassName}{% if non-KOMA class
  \IfFileExists{parskip.sty}{%
    \usepackage{parskip}
  }{% else
    \setlength{\parindent}{0pt}
    \setlength{\parskip}{6pt plus 2pt minus 1pt}}
}{% if KOMA class
  \KOMAoptions{parskip=half}}
\makeatother
\usepackage{xcolor}
\usepackage[margin=2.5cm]{geometry}
\usepackage{longtable,booktabs,array}
\usepackage{calc} % for calculating minipage widths
% Correct order of tables after \paragraph or \subparagraph
\usepackage{etoolbox}
\makeatletter
\patchcmd\longtable{\par}{\if@noskipsec\mbox{}\fi\par}{}{}
\makeatother
% Allow footnotes in longtable head/foot
\IfFileExists{footnotehyper.sty}{\usepackage{footnotehyper}}{\usepackage{footnote}}
\makesavenoteenv{longtable}
\setlength{\emergencystretch}{3em} % prevent overfull lines
\providecommand{\tightlist}{%
  \setlength{\itemsep}{0pt}\setlength{\parskip}{0pt}}
\setcounter{secnumdepth}{-\maxdimen} % remove section numbering
% ==========================================================================
%  IA na Sociedade — cabeçalho LaTeX (pdflatex)
%  Programa CiberExt 26-29 · FEELT38103
%  Adaptado de "AI in Society" (Universidade de Helsinque / Una Europa)
%  Licença: CC BY-NC-SA 4.0
% --------------------------------------------------------------------------
%  Espelha a identidade do ia-style.css: papel quente, tinta ardósia e a
%  faixa tricolor dos três eixos do curso (conceitos / aplicações / riscos).
% ==========================================================================

% ---------- Fontes ----------
% A serifada vem do build.sh, que passa  -V fontfamily=charter  ao pandoc.
% Aqui só acrescentamos a sem-serifa; ambas com fallback, porque nem toda
% instalação de LaTeX traz os mesmos pacotes.
\IfFileExists{helvet.sty}{\usepackage[scaled=0.90]{helvet}}{}
\usepackage{textcomp}
\IfFileExists{microtype.sty}{\usepackage{microtype}}{}

% ---------- Pacotes de layout ----------
\usepackage{xcolor}
\usepackage{tcolorbox}
\tcbuselibrary{skins,breakable}
\usepackage{titlesec}
\usepackage{fancyhdr}
\usepackage{enumitem}
\usepackage{amssymb}
\usepackage{booktabs}
\usepackage{array}
\usepackage{colortbl}
\usepackage{ragged2e}
\usepackage{fvextra}
\usepackage{newunicodechar}

% ---------- Caracteres Unicode que o pdflatex não conhece ----------
\newunicodechar{✓}{{\color{iaAplic}\checkmark}}
\newunicodechar{✅}{{\color{iaAplic}\checkmark}}
\newunicodechar{→}{\ensuremath{\rightarrow}}
\newunicodechar{←}{\ensuremath{\leftarrow}}
\newunicodechar{↔}{\ensuremath{\leftrightarrow}}
\newunicodechar{–}{--}
\newunicodechar{—}{---}
\newunicodechar{€}{\texteuro}
\newunicodechar{£}{\textsterling}
\newunicodechar{×}{\ensuremath{\times}}
\newunicodechar{≈}{\ensuremath{\approx}}
\newunicodechar{≥}{\ensuremath{\geq}}
\newunicodechar{≤}{\ensuremath{\leq}}
\newunicodechar{•}{\textbullet}
\newunicodechar{″}{''}
\newunicodechar{′}{'}

% ==========================================================================
%  PALETA — os mesmos valores do ia-style.css
% ==========================================================================
\definecolor{iaConceitos}{HTML}{24505C}   % petróleo — títulos, eixo 1
\definecolor{iaAplic}{HTML}{16746C}       % verdete  — subtítulos, eixo 2
\definecolor{iaRiscos}{HTML}{B45F14}      % açafrão  — destaques, eixo 3
\definecolor{iaAmeixa}{HTML}{77345F}      % links internos, termos técnicos
\definecolor{iaTinta}{HTML}{1F2933}       % corpo do texto
\definecolor{iaTintaFraca}{HTML}{5C6670}
\definecolor{iaRegra}{HTML}{E4DBCD}
\definecolor{iaBloco}{HTML}{F3EEE4}
\definecolor{iaNota}{HTML}{EAF1F0}
\definecolor{iaCaso}{HTML}{FCF3E6}
\definecolor{iaFilos}{HTML}{F7EFF4}

\color{iaTinta}

% ==========================================================================
%  ASSINATURA — faixa tricolor dos três eixos
% ==========================================================================
\newcommand{\faixaEixos}[1][\linewidth]{%
  \noindent
  \textcolor{iaConceitos}{\rule{0.3333#1}{4pt}}%
  \textcolor{iaAplic}{\rule{0.3333#1}{4pt}}%
  \textcolor{iaRiscos}{\rule{0.3333#1}{4pt}}%
  \par}

% ---------- Página de título ----------
\makeatletter
\newcommand{\subtitle}[1]{\gdef\@subtitle{#1}}
\providecommand{\@subtitle}{}
\def\@maketitle{%
  \newpage\null\vskip 1.5em%
  \faixaEixos
  \vskip 2.2em
  {\raggedright
    {\sffamily\bfseries\huge\color{iaConceitos}\@title\par}%
    \vskip 0.7em%
    {\sffamily\large\color{iaTintaFraca}\@subtitle\par}%
    \vskip 2em%
    {\color{iaRegra}\rule{\linewidth}{1pt}}%
    \vskip 1em%
    {\sffamily\small\color{iaTinta}\@author\par}%
    \vskip 0.5em%
    {\sffamily\footnotesize\color{iaTintaFraca}\@date\par}%
  }\par\vskip 2.5em}
\makeatother

% ==========================================================================
%  TÍTULOS DE SEÇÃO
% ==========================================================================
\titleformat{\section}[block]
  {\sffamily\Large\bfseries\color{iaConceitos}}
  {\thesection}{0.7em}{}
  [{\vspace{2pt}\color{iaRegra}\titlerule[1.2pt]}]
\titleformat{\subsection}
  {\sffamily\large\bfseries\color{iaAplic}}
  {\thesubsection}{0.6em}{}
\titleformat{\subsubsection}
  {\sffamily\normalsize\bfseries\color{iaRiscos}}
  {\thesubsubsection}{0.5em}{}
\titleformat{\paragraph}[runin]
  {\sffamily\small\bfseries\color{iaRiscos}}
  {}{0em}{}

\titlespacing*{\section}{0pt}{2.6ex plus 1ex minus .2ex}{1.3ex plus .2ex}
\titlespacing*{\subsection}{0pt}{2.2ex plus 1ex minus .2ex}{1ex plus .2ex}

% ==========================================================================
%  CABEÇALHO E RODAPÉ
% ==========================================================================
\setlength{\headheight}{24pt}
\renewcommand{\sectionmark}[1]{\markboth{#1}{}}
\pagestyle{fancy}
\fancyhf{}
\renewcommand{\headrulewidth}{0.6pt}
\renewcommand{\footrulewidth}{0pt}
\fancyhead[L]{\sffamily\footnotesize\color{iaAplic}Ética da IA}
\fancyhead[R]{\sffamily\footnotesize\color{iaAplic}\nouppercase{\leftmark}}
\fancyfoot[C]{\sffamily\footnotesize\color{iaConceitos}\thepage}
\renewcommand{\headrule}{%
  \hbox to\headwidth{\color{iaRegra}\leaders\hrule height \headrulewidth\hfill}}

% ==========================================================================
%  TEXTO CORRIDO
% ==========================================================================
\linespread{1.08}
\setlength{\parskip}{0.55em}
\setlength{\parindent}{0pt}
\emergencystretch=3em
\hbadness=10000

% ---------- Termos técnicos / código inline ----------
\let\oldtexttt\texttt
\renewcommand{\texttt}[1]{{\color{iaAmeixa}\small\ttfamily #1}}

% ---------- Blocos verbatim ----------
\fvset{breaklines=true,breakanywhere=true,fontsize=\small}

% ---------- Blocos de código (ambiente Shaded do pandoc) ----------
\renewenvironment{Shaded}
  {\begin{tcolorbox}[
      breakable, enhanced jigsaw, rounded corners,
      arc=3pt, boxrule=0pt, frame hidden,
      colback=iaBloco,
      borderline west={3pt}{0pt}{iaAmeixa},
      left=10pt, right=8pt, top=6pt, bottom=6pt,
      before skip=9pt, after skip=9pt]}
  {\end{tcolorbox}}

% ---------- Citações: caixa de caso ----------
\renewenvironment{quote}
  {\begin{tcolorbox}[
      breakable, enhanced jigsaw, rounded corners,
      arc=3pt, boxrule=0pt, frame hidden,
      colback=iaCaso,
      borderline west={3pt}{0pt}{iaRiscos},
      left=12pt, right=10pt, top=9pt, bottom=9pt,
      before skip=9pt, after skip=9pt]}
  {\end{tcolorbox}}

% ==========================================================================
%  BLOCOS DE APOIO
%  No markdown, escreva:   ::: nota   ...texto...   :::
%  O pandoc converte fenced_divs para os ambientes abaixo.
% ==========================================================================
\newtcolorbox{iabox}[3]{
  breakable, enhanced jigsaw, rounded corners,
  arc=3pt, boxrule=0pt, frame hidden,
  colback=#2,
  borderline west={3pt}{0pt}{#1},
  left=12pt, right=10pt, top=9pt, bottom=9pt,
  before skip=11pt, after skip=11pt,
  before upper={\setlength{\parskip}{0.5em}\setlength{\parindent}{0pt}},
  title={\sffamily\bfseries\footnotesize\textcolor{#1}{\MakeUppercase{#3}}},
  attach title to upper={\par\vspace{5pt}},
  fonttitle=\normalfont}

% Cada ambiente aceita um titulo opcional entre colchetes, vindo do
% atributo data-titulo do markdown:   \begin{filosofia}[Titulo] ... \end{filosofia}
% Os quatro primeiros espelham os icones do material original.
\newenvironment{filosofia}[1][Conceito filosofico]{\begin{iabox}{iaAmeixa}{iaFilos}{#1}}{\end{iabox}}
\newenvironment{tecnica}[1][Nota tecnica]{\begin{iabox}{iaConceitos}{iaNota}{#1}}{\end{iabox}}
\newenvironment{contexto}[1][Contexto]{\begin{iabox}{iaAplic}{iaBloco}{#1}}{\end{iabox}}
\newenvironment{definicao}[1][Definicao]{\begin{iabox}{iaRiscos}{iaCaso}{#1}}{\end{iabox}}

\newenvironment{nota}[1][Nota]{\begin{iabox}{iaConceitos}{iaNota}{#1}}{\end{iabox}}
\newenvironment{caso}[1][Caso]{\begin{iabox}{iaRiscos}{iaCaso}{#1}}{\end{iabox}}
\newenvironment{reflexao}[1][Para pensar]{\begin{iabox}{iaAplic}{iaBloco}{#1}}{\end{iabox}}
\newenvironment{licenca}
  {\par\vspace{2em}\faixaEixos\par\vspace{0.8em}
   \begingroup\sffamily\scriptsize\color{iaTintaFraca}}
  {\endgroup\par}

% ==========================================================================
%  LISTAS
% ==========================================================================
\setlist[itemize]{leftmargin=1.4em, itemsep=2pt, topsep=4pt, parsep=0pt}
\setlist[enumerate]{leftmargin=1.7em, itemsep=2pt, topsep=4pt, parsep=0pt}
\renewcommand{\labelitemi}{\textcolor{iaRiscos}{\textbullet}}
\renewcommand{\labelitemii}{\textcolor{iaAplic}{--}}

% ==========================================================================
%  TABELAS
% ==========================================================================
\renewcommand{\arraystretch}{1.3}
\arrayrulecolor{iaRegra}

% ==========================================================================
%  RÓTULOS EM PORTUGUÊS
%  Definidos aqui (e não via -V na linha de comando) porque acentos passados
%  como argumento dependem do locale do terminal e podem se corromper.
% ==========================================================================
\renewcommand{\contentsname}{Sumário}
\renewcommand{\refname}{Referências}
\renewcommand{\figurename}{Figura}
\renewcommand{\tablename}{Tabela}
\ifLuaTeX
  \usepackage{selnolig}  % disable illegal ligatures
\fi
\IfFileExists{bookmark.sty}{\usepackage{bookmark}}{\usepackage{hyperref}}
\IfFileExists{xurl.sty}{\usepackage{xurl}}{} % add URL line breaks if available
\urlstyle{same}
\hypersetup{
  pdftitle={Ética da Inteligência Artificial},
  colorlinks=true,
  linkcolor={iaAmeixa},
  filecolor={Maroon},
  citecolor={Blue},
  urlcolor={iaRiscos},
  pdfcreator={LaTeX via pandoc}}

\title{Ética da Inteligência Artificial}
\usepackage{etoolbox}
\makeatletter
\providecommand{\subtitle}[1]{% add subtitle to \maketitle
  \apptocmd{\@title}{\par {\large #1 \par}}{}{}
}
\makeatother
\subtitle{Capítulo 1 --- Exercícios}
\author{Tradução e adaptação para o português: José Lucas Lira Bizil,
Fernando Mazzeto Lisboa Lima e Matheus da Silva Fernandes\\
Programa CiberExt 26-29 · FEELT38103 · Universidade Federal de
Uberlândia\\
Original: \emph{Ethics of AI}, Universidade de Helsinque}
\date{2026}

\begin{document}
\maketitle

{
\hypersetup{linkcolor=iaConceitos}
\setcounter{tocdepth}{2}
\tableofcontents
}
\hypertarget{capuxedtulo-1-exercuxedcios}{%
\section{Capítulo 1 --- Exercícios}\label{capuxedtulo-1-exercuxedcios}}

\begin{nota}[Sobre estes exercícios]

Traduzidos dos questionários do curso original. Ao lado de cada um está
indicado o componente correspondente na Open edX.

\end{nota}

\hypertarget{exercuxedcio-1a-valores-e-normas-muxfaltipla-escolha-seleuxe7uxe3o-muxfaltipla}{%
\subsection{\texorpdfstring{Exercício 1a --- Valores e normas
\texttt{{[}múltipla\ escolha\ ·\ seleção\ múltipla{]}}}{Exercício 1a --- Valores e normas {[}múltipla escolha · seleção múltipla{]}}}\label{exercuxedcio-1a-valores-e-normas-muxfaltipla-escolha-seleuxe7uxe3o-muxfaltipla}}

\textbf{Pontuação:} 1 ponto · \textbf{Tentativas:} ilimitadas
\textbf{Componente Open edX:} \texttt{choiceresponse} (caixas de
seleção)

Todos nós temos nossos próprios valores pessoais. Avalie a importância
dos valores a seguir para você, \textbf{selecionando até 5 valores mais
importantes}. (Adaptado de Schwartz, 2001.)

\begin{itemize}
\tightlist
\item
  Honrar os pais e os mais velhos
\item
  Autodisciplina
\item
  Polidez
\item
  Evitar violar expectativas ou normas sociais
\item
  Respeito às outras pessoas
\item
  Entusiasmo, novidade e desafio na vida
\item
  Prazer
\item
  Coragem
\item
  Proteção do bem-estar de todas as pessoas
\item
  Igualdade entre todas as pessoas
\item
  Criatividade
\item
  Independência de pensamento e de escolha da ação
\item
  Lealdade
\item
  Controle ou domínio sobre pessoas e recursos
\item
  Ordem social
\item
  Segurança
\item
  Honestidade
\item
  Proteção da natureza
\item
  Status social e prestígio
\item
  Limpeza
\item
  Estabilidade da sociedade
\end{itemize}

\begin{center}\rule{0.5\linewidth}{0.5pt}\end{center}

\hypertarget{exercuxedcio-1b-valores-e-normas-escala-de-avaliauxe7uxe3o}{%
\subsection{\texorpdfstring{Exercício 1b --- Valores e normas
\texttt{{[}escala\ de\ avaliação{]}}}{Exercício 1b --- Valores e normas {[}escala de avaliação{]}}}\label{exercuxedcio-1b-valores-e-normas-escala-de-avaliauxe7uxe3o}}

\textbf{Pontuação:} 1 ponto · \textbf{Tentativas:} ilimitadas
\textbf{Componente Open edX:} \texttt{multiplechoiceresponse} (uma
escala por item) ou componente de pesquisa (\emph{survey})

Suponha que você seja nomeado membro de um grupo de especialistas em IA
do governo. A tarefa do grupo é decidir os códigos éticos para a
implantação da IA em seu país. Pede-se que você selecione os valores que
orientarão e governarão o uso da IA nos próximos anos, votando numa
consulta.

Marque a importância de cada valor da lista a seguir, numa escala de
\textbf{1 (nada importante) a 7 (extremamente importante)}:

\begin{longtable}[]{@{}ll@{}}
\toprule\noalign{}
Valor & Escala \\
\midrule\noalign{}
\endhead
\bottomrule\noalign{}
\endlastfoot
Status social e prestígio & 1--7 \\
Controle ou domínio sobre pessoas e recursos & 1--7 \\
Prazer & 1--7 \\
Entusiasmo, novidade e desafio na vida & 1--7 \\
Individualidade & 1--7 \\
Independência de pensamento e de escolha da ação & 1--7 \\
Criatividade & 1--7 \\
Igualdade entre todas as pessoas & 1--7 \\
Proteção do bem-estar de todas as pessoas & 1--7 \\
Proteção da natureza & 1--7 \\
Honestidade & 1--7 \\
Lealdade & 1--7 \\
Polidez & 1--7 \\
Autodisciplina & 1--7 \\
Honrar os pais e os mais velhos & 1--7 \\
Estabilidade da sociedade & 1--7 \\
Segurança & 1--7 \\
Limpeza & 1--7 \\
\end{longtable}

\begin{reflexao}[Observação didática]

Os exercícios 1a e 1b não têm resposta certa: servem para o estudante
perceber a diferença entre os valores que sustenta pessoalmente e
aqueles que consideraria adequados para governar uma política pública.
Vale explicitar isso no ambiente do curso, para que ninguém tente
``acertar''.

\end{reflexao}

\begin{center}\rule{0.5\linewidth}{0.5pt}\end{center}

\hypertarget{exercuxedcio-2-valores-e-normas-associauxe7uxe3o}{%
\subsection{\texorpdfstring{Exercício 2 --- Valores e normas
\texttt{{[}associação{]}}}{Exercício 2 --- Valores e normas {[}associação{]}}}\label{exercuxedcio-2-valores-e-normas-associauxe7uxe3o}}

\textbf{Pontuação:} 1 ponto · \textbf{Tentativas:} 1 \textbf{Componente
Open edX:} \texttt{optionresponse} (menu suspenso por item)

Associe cada norma abaixo ao valor correspondente.

\begin{enumerate}
\def\labelenumi{\arabic{enumi}.}
\item
  ``Proteções adicionais devem ser dadas às pessoas cujos dados foram
  usados para treinar modelos de IA, como o direito de acessar os
  modelos; de saber de onde se originaram e a quem estão sendo
  comercializados ou transmitidos; o direito de apagar a si mesmo de um
  modelo treinado; e o direito de expressar o desejo de que o modelo não
  seja usado no futuro.'' \emph{(Parlamento Europeu: The ethics of
  artificial intelligence: Issues and initiatives)}
\item
  ``O desenvolvimento da IA deve promover a justiça e buscar eliminar
  todos os tipos de discriminação.'' \emph{(AI4People: An Ethical
  Framework for a Good AI Society)}
\item
  ``A IA deve promover o florescimento humano individual e ampliar a
  agência intelectual, libertando as pessoas do trabalho rotineiro.''
  \emph{(AI4People: An Ethical Framework for a Good AI Society)}
\item
  A tecnologia baseada em IA não deve ser usada sistematicamente para
  manipular cidadãos em eleições, prejudicando a democracia.
\item
  A pegada de carbono da IA deve ser reduzida.
\item
  ``Sistemas autônomos não devem prejudicar a liberdade dos seres
  humanos de estabelecer seus próprios padrões e normas intelectuais.''
  \emph{(AI4People: An Ethical Framework for a Good AI Society)}
\item
  ``Sistemas de IA devem ser tecnicamente robustos, e deve-se assegurar
  que não estejam abertos a uso malicioso.'' \emph{(Comissão Europeia:
  Ethics Guidelines for Trustworthy AI)}
\end{enumerate}

\textbf{Opções disponíveis} (cada uma usada uma vez):

\begin{itemize}
\tightlist
\item
  Privacidade
\item
  Igualdade e não discriminação
\item
  Autorrealização e florescimento humano
\item
  Democracia
\item
  Proteção da natureza
\item
  Autonomia
\item
  Segurança
\end{itemize}

\begin{center}\rule{0.5\linewidth}{0.5pt}\end{center}

\hypertarget{exercuxedcio-3-tuitando-sobre-questuxf5es-uxe9ticas-dissertativa}{%
\subsection{\texorpdfstring{Exercício 3 --- Tuitando sobre questões
éticas
\texttt{{[}dissertativa{]}}}{Exercício 3 --- Tuitando sobre questões éticas {[}dissertativa{]}}}\label{exercuxedcio-3-tuitando-sobre-questuxf5es-uxe9ticas-dissertativa}}

\textbf{Pontuação:} 1 ponto · \textbf{Extensão:} 20 a 1000 palavras ·
\textbf{Tentativas:} ilimitadas \textbf{Componente Open edX:}
\texttt{openassessment} (revisão por pares)

Imagine que um dia você acabe numa discussão acalorada no Twitter. Ela
começa com o tuíte de um professor universitário (@TuringLives) sobre um
modelo de tradução de imagens. O modelo transformou uma imagem de
entrada pixelada da primeira-ministra finlandesa Sanna Marin na foto de
um homem branco de meia-idade.

\textbf{Observação:} este não é um caso real de recriação fotográfica
por IA.

\emph{Imagem 1: CC BY 4.0 Laura Kotila / Gabinete da Primeira-Ministra
da Finlândia (editada a partir do original). Imagens 2 e 3 são impressão
artística. Imagem 4: CC BY-NC 4.0 NVIDIA Corporation.}

\textbf{O que são algoritmos de tradução de imagens?}

Muitos dos exemplos mais conhecidos de tradução de imagens são
produzidos por uma rede generativa adversarial (\emph{Generative
Adversarial Network}, GAN). Uma GAN é um tipo de arquitetura de rede
neural voltada à modelagem generativa.

Nas GANs, duas redes competem entre si. Uma é treinada para gerar, por
exemplo, imagens semelhantes às dos dados de treinamento (gatos, rostos
humanos ou outras coisas). A tarefa da outra (chamada rede adversarial)
é separar as imagens geradas pela primeira das imagens reais dos dados
de treinamento.

O sistema treina os dois modelos lado a lado. Na primeira fase, a tarefa
do modelo adversarial é distinguir as imagens reais das tentativas
desajeitadas do modelo generativo. Contudo, à medida que a rede
generativa vai melhorando, o modelo adversarial também precisa melhorar.
O ciclo continua até que, por fim, as imagens geradas sejam quase
indistinguíveis das reais.

As GANs não tentam apenas reproduzir os itens dos dados de treinamento:
o sistema é treinado de modo a ter de gerar itens novos e de aparência
real. No entanto, as GANs --- como muitos outros algoritmos
contemporâneos --- produzem resultados que refletem os padrões
estatísticos dos dados de entrada.

Essas imagens foram desenvolvidas num projeto de pesquisa de Tero
Karras, Samuli Laine, Timo Aila e Jaakko Lehtinen no NVIDIA Research
Helsinki.

As GANs podem ser usadas para muitos fins, como tarefas de tradução de
imagem para imagem --- converter fotos de noite em fotos de dia --- ou
gerar fotos fotorrealistas de itens, objetos, cenas ou pessoas.

Para ver como as GANs funcionam, acesse:
\url{https://developer.ibm.com/articles/generative-adversarial-networks-explained/}

\textbf{Sua tarefa:} entre na discussão do Twitter e responda aos tuítes
formulando seu próprio ponto de vista sobre o assunto. Crie um nome de
usuário para si mesmo e escreva sua resposta.

\begin{itemize}
\tightlist
\item
  \textbf{Nome no Twitter:} \_\_\_\_\_\_\_\_\_\_\_\_\_\_\_
\item
  \textbf{Resposta:} \_\_\_\_\_\_\_\_\_\_\_\_\_\_\_
\end{itemize}

\begin{center}\rule{0.5\linewidth}{0.5pt}\end{center}

\hypertarget{exercuxedcio-3c-tuitando-sobre-questuxf5es-uxe9ticas-dissertativa}{%
\subsection{\texorpdfstring{Exercício 3c --- Tuitando sobre questões
éticas
\texttt{{[}dissertativa{]}}}{Exercício 3c --- Tuitando sobre questões éticas {[}dissertativa{]}}}\label{exercuxedcio-3c-tuitando-sobre-questuxf5es-uxe9ticas-dissertativa}}

\textbf{Pontuação:} 3 pontos · \textbf{Extensão:} 150 a 300 palavras ·
\textbf{Tentativas:} ilimitadas \textbf{Componente Open edX:}
\texttt{openassessment} (revisão por pares)

Sobre o que trata a discussão do Twitter acima --- ou seja, quais são os
valores, normas e princípios éticos em jogo? Você os considera
significativos? Analise a discussão e explique os temas éticos. Escreva
um ensaio curto (150 a 300 palavras) em que você julgue a importância
desses tópicos.

Leia as instruções com atenção. Sua resposta será revisada por outros
participantes e pelos instrutores. Confira sua resposta antes de enviar:
uma vez enviada, não poderá mais editá-la.

\begin{center}\rule{0.5\linewidth}{0.5pt}\end{center}

\hypertarget{avaliauxe7uxe3o-do-capuxedtulo-pesquisa}{%
\subsection{\texorpdfstring{Avaliação do capítulo
\texttt{{[}pesquisa{]}}}{Avaliação do capítulo {[}pesquisa{]}}}\label{avaliauxe7uxe3o-do-capuxedtulo-pesquisa}}

\textbf{Pontuação:} 1 ponto · \textbf{Tentativas:} ilimitadas

Hora de dar retorno. O que você achou deste capítulo?

\begin{itemize}
\tightlist
\item
  \textbf{Utilidade deste capítulo} --- nota de 1 a 5 (1 = pouco útil /
  5 = muito útil)
\item
  \textbf{Clareza deste capítulo} --- nota de 1 a 5 (1 = pouca clareza /
  5 = boa clareza)
\item
  \textbf{Dificuldade deste capítulo} --- nota de 1 a 5 (1 = nada
  difícil / 5 = muito difícil)
\item
  \textbf{Outros comentários} --- até 600 palavras
\end{itemize}

\begin{center}\rule{0.5\linewidth}{0.5pt}\end{center}

\begin{licenca}

\textbf{Material original:} \emph{Ethics of AI}, um curso online
gratuito da Universidade de Helsinque. Professores responsáveis:
Anna-Mari Rusanen e Jukka K. Nurminen. Demais contribuintes principais:
Santeri Räisänen, Sasu Tarkoma e Saara Halmetoja. Disponível em
\url{https://ethics-of-ai.mooc.fi/}.

\textbf{Esta versão:} tradução e adaptação para o português brasileiro
realizada por José Lucas Lira Bizil, Fernando Mazzeto Lisboa Lima e
Matheus da Silva Fernandes, no âmbito do Programa CiberExt 26-29
(FEELT38103), Universidade Federal de Uberlândia, 2026. Obra derivada;
os autores originais não a revisaram nem a endossam.

\textbf{Licença:} CC BY-NC-SA 4.0 ---
\url{https://creativecommons.org/licenses/by-nc-sa/4.0/deed.pt-br}.

\end{licenca}

\end{document}
